\section{Secondary Endpoints and Multiplicity Adjustment}

\subsection{Hierarchical Testing Procedure}

All confirmatory secondary endpoints will be tested using the fixed-sequence (hierarchical)
procedure outlined below. Testing proceeds to the next endpoint only if the preceding
test is significant at $\alpha = 0.025$ (one-sided).

\subsubsection{Multiplicity Adjustment Flowchart}

\begin{center}
\begin{tikzpicture}[
  node distance=7mm and 4mm,
  box/.style={
    rectangle, rounded corners=3pt,
    draw=sapblue, fill=sapfog, thick,
    text width=72mm, align=center,
    inner sep=5pt, font=\small
  },
  decision/.style={
    diamond, draw=sapsteel, fill=white, thick,
    aspect=2.8, align=center, font=\footnotesize,
    inner sep=2pt
  },
  stop/.style={
    rectangle, rounded corners=3pt,
    draw=sapwarn, fill=sapwarnbg, thick,
    text width=52mm, align=center,
    inner sep=4pt, font=\small\bfseries\color{sapwarn}
  },
  arr/.style={-Stealth, thick, color=sapblue},
  nos/.style={-Stealth, thick, color=sapwarn}
]

\node[box] (ep1) {\textbf{H1:} Primary — Trough FEV\textsubscript{1} at Wk 52\\\small(MMRM, one-sided $\alpha=0.025$)};

\node[decision, below=of ep1] (d1) {Sig?};

\node[box, below=of d1] (ep2) {\textbf{H2:} SGRQ Total Score at Wk 52\\\small(MMRM)};

\node[decision, below=of ep2] (d2) {Sig?};

\node[box, below=of d2] (ep3) {\textbf{H3:} Exacerbation Rate over 52 Wks\\\small(Negative Binomial)};

\node[decision, below=of ep3] (d3) {Sig?};

\node[box, below=of d3] (ep4) {\textbf{H4:} Peak FEV\textsubscript{1} at Wk 52\\\small(MMRM)};

\node[stop, right=24mm of d1] (stop1) {Stop confirmatory testing};
\node[stop, right=24mm of d2] (stop2) {Stop confirmatory testing};
\node[stop, right=24mm of d3] (stop3) {Stop confirmatory testing};

\draw[arr] (ep1) -- (d1);
\draw[arr] (d1) -- node[left,font=\footnotesize]{Yes} (ep2);
\draw[arr] (ep2) -- (d2);
\draw[arr] (d2) -- node[left,font=\footnotesize]{Yes} (ep3);
\draw[arr] (ep3) -- (d3);
\draw[arr] (d3) -- node[left,font=\footnotesize]{Yes} (ep4);

\draw[nos] (d1) -- node[above,font=\footnotesize]{No} (stop1);
\draw[nos] (d2) -- node[above,font=\footnotesize]{No} (stop2);
\draw[nos] (d3) -- node[above,font=\footnotesize]{No} (stop3);

\end{tikzpicture}
\end{center}

\begin{tcolorbox}[sapnote, title={\faIcon{shield-alt}\ Family-Wise Error Rate Control}]
  The fixed-sequence procedure controls the family-wise error rate (FWER) at one-sided
  $\alpha = 0.025$ under arbitrary correlation among test statistics. No Bonferroni,
  Holm, or Hochberg adjustments are applied within the sequence.
  Exploratory endpoints (H5 onward) are not multiplicity-adjusted; $p$-values are nominal.
\end{tcolorbox}

\subsection{Secondary Analysis Methods}

\begin{tabularx}{\linewidth}{@{}
  >{\centering\arraybackslash}p{8mm}
  >{\raggedright\arraybackslash}X
  >{\raggedright\arraybackslash}p{38mm}
  >{\centering\arraybackslash}p{16mm}@{}}
\toprule
\rowcolor{sapfog}
\textbf{H\#} & \textbf{Endpoint} & \textbf{Method} & \textbf{TFL Ref} \\
\midrule
H2 & SGRQ Total Score change from baseline at Wk 52 & MMRM (same model) & ET-EFF-002 \\
\addlinespace[2pt]
H3 & Annualised rate of moderate/severe COPD exacerbations & Negative Binomial GLM, log(time on study) offset & ET-EFF-003 \\
\addlinespace[2pt]
H4 & Peak FEV\textsubscript{1} (2h post-dose) change from BL at Wk 52 & MMRM & ET-EFF-004 \\
\addlinespace[2pt]
H5 & Proportion of FEV\textsubscript{1} responders ($\geq$100 mL) at Wk 52 & Logistic regression; OR + 95\% CI & ET-EFF-005 \\
\addlinespace[2pt]
H6 & Time to first moderate/severe exacerbation & Kaplan--Meier; log-rank; Cox PH & ET-EFF-006 \\
\bottomrule
\end{tabularx}

\clearpage

% ============================================================
%  SECTION 6 — SAFETY ANALYSES
% ============================================================
